\chapter{Documentation}
\label{ch:documentation}

The previous chapter was about metadata you \emph{add}. This one is about
metadata you already have.

Your library is probably documented. Its headers carry \code{///} lines and
\code{/** @param \ldots */} blocks, its declarations name their parameters, and
its defaults are written down where they belong --- in the C++. Until you bind
it, none of that reaches the people calling your library from Python or
TypeScript: an annotation-only pipeline can say nothing your headers did not
repeat, in rosetta's own vocabulary, a second time.

rosetta reads the headers it is already opening. The result is that a generated
binding arrives documented, with keyword arguments and working defaults, and
you write nothing.

\section{Turning it on}

There is nothing to turn on. Harvesting is the default; the only setting is how
to turn it \textbf{off}.

\begin{lstlisting}[style=json]
{ "doc_comments": false }     // default: true
\end{lstlisting}

Set it to \code{false} to generate exactly what earlier versions of rosetta did
--- documentation from annotations alone, positional parameter names, no
defaults. There is no per-class or per-target switch: the harvest is a property
of the \emph{manifest}, and the backends decide what to do with what it finds.

\section{An example, end to end}

A header that never mentions rosetta:

\begin{lstlisting}[style=cpp]
/// The direction a polygon's vertices are wound in.
enum class Winding {
    CCW, ///< Counter-clockwise: a positive signed area.
    CW,  ///< Clockwise: a negative signed area.
};

/**
 * @brief A closed polygon, stored as parallel coordinate arrays.
 */
class Polygon {
public:
    /// The x coordinate of each vertex, in order.
    std::vector<double> xs;

    /**
     * @brief Resample the outline at evenly spaced points.
     *
     * @param samples how many points to produce
     * @param closed  whether to repeat the first point at the end
     * @return the coordinates, x and y interleaved
     */
    std::vector<double> resample(int samples = 64, bool closed = true) const;
};
\end{lstlisting}

\noindent
becomes, in Python:

\begin{lstlisting}[style=py]
>>> help(polygon.Polygon.resample)
resample(self, samples: int = 64, closed: bool = True) -> list[float]

    Resample the outline at evenly spaced points.

    Args:
        samples: how many points to produce
        closed: whether to repeat the first point at the end

    Returns:
        the coordinates, x and y interleaved

>>> p.resample(samples=8, closed=False)   # keyword arguments
>>> p.resample()                          # the C++ defaults apply
\end{lstlisting}

\noindent
and, in the \code{.d.ts}:

\begin{lstlisting}[style=js]
/**
 * Resample the outline at evenly spaced points.
 * @param samples how many points to produce
 * @param closed whether to repeat the first point at the end
 * @returns the coordinates, x and y interleaved
 */
resample(samples?: number, closed?: boolean): number[];
\end{lstlisting}

\section{What is read}

\begin{center}
\small
\begin{tabular}{@{}L{5.4cm}L{8.4cm}@{}}
\toprule
\textbf{In the header} & \textbf{In the IR} \\
\midrule
\code{///}, \code{//!}, \code{/** */}, \code{/*! */}, and trailing \code{///<} &
  the member's documentation (\code{GenField::doc}, \code{GenMethod::doc}) \\
\addlinespace
\code{@brief}, \code{@details}, untagged prose & the same \\
\addlinespace
\code{@param <name> \ldots}, including the \code{[in]} / \code{[out]} /
  \code{[in,out]} markers and text running over several lines &
  \code{GenParam::doc} --- per argument, matched \textbf{by name} \\
\addlinespace
\code{@return}, \code{@returns}, \code{@result} & \code{GenMethod::returns} \\
\addlinespace
\code{@note}, \code{@warning}, \code{@deprecated}, \code{@throws},
  \code{@pre}, \code{@post}, \code{@see}, \code{@since} &
  labelled paragraphs appended to the documentation \\
\addlinespace
a comment above a \code{class} / \code{struct} & \code{GenClass::brief} \\
\addlinespace
\textbf{parameter names} in the declaration & \code{GenParam::name} \\
\addlinespace
\textbf{default arguments} (\code{int samples = 64}) &
  \code{GenParam::has\_default} and \code{GenParam::default\_text} \\
\bottomrule
\end{tabular}
\end{center}

\code{@tparam}, \code{@file}, \code{@ingroup}, \code{@defgroup} and the other
bookkeeping commands are dropped --- they describe the source, not the API.

\section{Which backend does what with it}

Every backend receives the same IR. What each one can \emph{do} with it differs,
because the host languages differ.

\begin{center}
\small
\begin{tabular}{@{}lL{10.4cm}@{}}
\toprule
\textbf{Target} & \textbf{What it emits} \\
\midrule
\lang{python} \lang{nanobind} &
  The docstring, with \code{Args:} and \code{Returns:} sections folded in ---
  Python has one place for all of it. Parameter names become
  \code{py::arg("\ldots")} / \code{nb::arg("\ldots")}, so calls take keyword
  arguments; recovered defaults become real Python defaults. The class comment
  is the type's docstring. \\
\addlinespace
\lang{typescript} &
  A TSDoc block: the prose, one \code{@param} per documented argument, and
  \code{@returns}. A defaulted parameter is declared \textbf{optional}
  (\code{samples?: number}), so an editor offers the short call. \\
\addlinespace
\lang{markdown} \lang{html} &
  The signature carries names and defaults; the \code{@param} text becomes a
  list under each method, \code{@return} its own line. The class comment opens
  the section. \\
\addlinespace
\lang{openapi} &
  The operation \code{description}; each positional argument gets a
  \code{title} (its name), a \code{description} (its \code{@param} text) and a
  \code{default}. The wire shape does not change. \\
\addlinespace
\lang{csharp} &
  A \code{/// <summary>} XML doc block, and real parameter names in the wrapper
  signature. \\
\addlinespace
\lang{java} &
  A Javadoc block, and real parameter names in the wrapper signature. \\
\addlinespace
\lang{qt} \lang{qml} \lang{imgui} &
  A field's documentation becomes its tooltip, and parameter names label the
  argument boxes of a method's call row --- an inspector built from a header
  that was never written for one. \\
\addlinespace
\lang{node} \lang{wasm} \lang{lua} \lang{julia} &
  Nothing. These runtimes have no docstring slot and their emitted wrappers
  spell positional locals, so the names and text sit unused in the IR. For
  \lang{node} and \lang{wasm}, \lang{typescript} is where they land ---
  which is the argument for adding it. \\
\addlinespace
\lang{dynamic} &
  Names and text sit in the emitted \code{MetaClass} tables, so a UI built by
  querying the registry can show them without ever including your headers. \\
\bottomrule
\end{tabular}
\end{center}

\begin{note}
Add \lang{typescript} to a manifest that already has \lang{node} or
\lang{wasm}. It costs one line, compiles nothing, and it is the only place
those two targets' documentation can go.
\end{note}

\section{Default arguments, and why they are different}

Parameter \emph{names} come from reflection: \code{std::meta::identifier\_of}
reads them from the declaration. So does the \emph{fact} of a default ---
\code{std::meta::has\_default\_argument}. But P2996 offers no way to ask what
the default \emph{expression was}, so its spelling is read from the header text
instead.

The two readings are of the same declaration, and they are
\textbf{cross-checked}: a spelling is used only where reflection independently
confirms there is a default to spell. Where they disagree, the text loses.

That leaves the question of which expressions rosetta is willing to repeat into
generated C++. The bar is high, because a wrong answer here does not produce a
poor docstring --- it produces a binding that does not compile.

\begin{center}
\small
\begin{tabular}{@{}L{6.6cm}L{7.2cm}@{}}
\toprule
\textbf{Accepted} & \textbf{Refused} \\
\midrule
literals: \code{64}, \code{1e-6}, \code{true}, \code{'x'},
  \code{"polygon"} &
  anything with a call: \code{Point(0, 0)}, \code{sizeof(int)} \\
\addlinespace
\code{\{\}} &
  an expression: \code{a + b}, \code{n * 2} \\
\addlinespace
qualified-ids: \code{Winding::CCW},
  \code{Limits::Max} &
  a bare identifier: \code{kDefaultEpsilon} --- it would not resolve where the
  binding spells it \\
\bottomrule
\end{tabular}
\end{center}

A parameter whose default is refused is still \emph{known} to be optional, and
the backends that only need that say so: \code{eps?: number} in the
\code{.d.ts}, \code{required: false} in OpenAPI, \code{eps?} in the Markdown
signature. Only Python, which needs the value itself, falls back to a required
argument.

\begin{warn}
The Python family adds a rule of its own. pybind and nanobind convert an
\code{arg} default into a Python object at \textbf{module import}, before any
call --- so a default whose type is not registered yet does not degrade
gracefully, it raises \code{arg(): could not convert default argument} and the
whole module fails to import. Defaults are therefore emitted only for numbers,
booleans, strings, and \textbf{enums bound in the same module}. Bind the enum
(list it in \field{classes}) and \code{Winding::CCW} becomes a working Python
default; leave it out and the default is silently withheld rather than breaking
the import.
\end{warn}

\section{Precedence}

Harvested text only ever fills a blank. In order:

\begin{enumerate}[leftmargin=*]
  \item an inline \code{[[= rosetta::doc\{"\ldots"\}]]} annotation --- the
    harvester skips a member that carries one outright;
  \item a JSON side-car entry (chapter~\ref{ch:annotations}), or a
    \field{doc} on a manifest \field{functions} entry;
  \item the header's own comment.
\end{enumerate}

So adopting this cannot change a binding you have already annotated, and the
two mechanisms compose: annotate the handful of members that need a
\code{range}, a \code{readonly} or a rewritten description, and let the headers
document everything else.

\section{What it will not do}

The harvester is a \textbf{scanner, not a compiler}. It tracks namespaces, class
bodies, access specifiers, comments, string and character literals and
preprocessor lines, and it reads declarations well enough to get a name, its
parameters and its defaults --- but it does not expand macros and does not
follow \code{\#include}.

Which is why nothing it finds is used on sight. Every entry is matched against
the \textbf{reflected} signature before it is applied, by arity and by
parameter names, and dropped when the two disagree. The consequences are worth
stating plainly:

\begin{itemize}[leftmargin=*]
  \item \textbf{Ambiguous overloads document neither.} Two entries of the same
    arity whose names cannot tell them apart are both discarded --- a plausible
    sentence on the wrong overload is worse than no sentence.
  \item \textbf{A stale comment loses.} A \code{@param radius} naming an
    argument the signature no longer has is dropped, not placed by position.
  \item \textbf{Private members are never read.} A documented private field is
    invisible to every backend, as it should be.
  \item \textbf{Constructors, destructors and operators carry no documentation}
    --- the IR has no slot for it. Constructor parameter \emph{names} do reach
    \code{py::init} and the TypeScript \code{constructor(\ldots)}.
  \item \textbf{Inherited members need their base bound.} A member flattened in
    from a base declared in a header no \field{classes} entry names is not
    covered; list that base as a class to reach it.
  \item \textbf{Template classes are not matched.} The IR names a
    specialisation (\code{SpatialFieldND<2>}); the header declares a template.
\end{itemize}

\begin{note}
Harvesting reads only the headers the manifest names --- the ones listed by
\field{classes}, their \field{extensions}, and \field{functions} --- resolved
against \field{user\_include} exactly as the compiler will resolve them. A
header that resolves nowhere contributes nothing and is not an error: it may be
a \field{generated\_headers} entry that does not exist yet.
\end{note}

\section{Where to look next}

\code{examples/doxygen} is this chapter as a runnable project: a documented
header, a three-line manifest, and a \code{test.py} that asserts against the
compiled module --- including all three of the cases above where rosetta
deliberately does less than you might expect. The manifest field is in
appendix~\ref{app:manifest}; the reference document is
\code{docs/MANIFEST.md}.
